\section{Measurement noise and the latent RNA correlation}\label{sec:noise}
Population means average over cells, and hence over the sampling noise of
single-cell RNA counts; a recipient's single-cell covariance does not. The
next result states what this does to Proposition~\ref{prop:means}.

\begin{proposition}[Latent and measured RNA]\label{prop:noise}
In every population let each cell have a latent RNA state $z$ and measured RNA
$x=z+e$, with $E(e\mid z,y)=0$ and $\operatorname{Cov}(e\mid z)=N(z)$
diagonal, and let $E(y\mid z,s)=Az+b$ with $A$ shared by all populations.
Write $\Sigma_z$ and $\Sigma_x$ for a recipient's latent and measured RNA
covariances and $\bar N=E\,N(z)$.
\begin{enumerate}
\item Population means satisfy $E(y\mid s)=A\,E(x\mid s)+b$, so they identify
$A$, the channel acting on the latent state, on $M$
(Proposition~\ref{prop:means}(1)).
\item $\Sigma_x=\Sigma_z+\bar N$ and $\operatorname{Cov}(x,y)=\Sigma_zA^\top$;
the cell-level regression slope of $y$ on $x$ is
$A\Sigma_z\Sigma_x^{-1}$, which differs from $A$ whenever $\bar N\neq0$.
\item Let $F(\Sigma)=Q_N^\top\Sigma^{1/2}A^\top\Sigma_y^{-1/2}$ with $Q_N$ an
orthonormal basis of $\Sigma^{-1/2}M$. The cross-covariances compatible with
$AP_M$ and with a valid joint law of $(z,y)$ exist if and only if
$\|F(\Sigma_z)\|_{\mathrm{op}}\le1$, and
$\|F(\Sigma_x)\|_{\mathrm{op}}\ge\|F(\Sigma_z)\|_{\mathrm{op}}$. Measurement
noise can therefore empty the set of Proposition~\ref{prop:means}(4) computed
with $\Sigma_x$ when the latent set is not empty.
\item If $x$ is a function of counts $c\sim\mathrm{Poisson}(\lambda)$ given
$\lambda$, the binomial halves $c_1\mid c\sim\mathrm{Bin}(c,1/2)$ and
$c_2=c-c_1$ are independent given $\lambda$, so any functions $f(c_1)$ and
$f(c_2)$ have covariance $\operatorname{Var}\{E(f(c_1)\mid\lambda)\}$, the
latent variance at half depth, which requires no paired cell.
\end{enumerate}
\end{proposition}

\begin{proof}
(1) $E(y\mid s)=E\{E(y\mid z,s)\mid s\}=A\,E(z\mid s)+b$ and
$E(x\mid s)=E(z\mid s)$. (2) By the law of total covariance
$\Sigma_x=\Sigma_z+E\,N(z)$; $\operatorname{Cov}(e,y)=0$ because
$E(e\mid z,y)=0$, so $\operatorname{Cov}(x,y)=\operatorname{Cov}(z,y)
=\Sigma_zA^\top$. (3) The first statement is Proposition~\ref{prop:means}(4)
applied to $(z,y)$. With $V$ an orthonormal basis of $M$ and $A_M=AV$,
$F(\Sigma)^\top F(\Sigma)=\Sigma_y^{-1/2}A_M(V^\top\Sigma^{-1}V)^{-1}A_M^\top\Sigma_y^{-1/2}$,
because $\Sigma^{1/2}Q_NQ_N^\top\Sigma^{1/2}=V(V^\top\Sigma^{-1}V)^{-1}V^\top$.
$\Sigma_x\succeq\Sigma_z$ gives $V^\top\Sigma_x^{-1}V\preceq V^\top\Sigma_z^{-1}V$
and, inverting, $F(\Sigma_x)^\top F(\Sigma_x)\succeq F(\Sigma_z)^\top F(\Sigma_z)$.
(4) Binomial thinning of a Poisson count gives independent Poisson counts with
half the mean; functions of independent variables are independent, and the
law of total covariance gives the stated covariance.
\end{proof}

In the analyses, $x$ is log-normalized expression. For each gene, the
correlation across cells of the log-normalized values of the two halves is its
split-half reliability; the Spearman--Brown formula $2r/(1+r)$ gives the
reliability at full depth, and one minus this is the share $\nu_j$ of the
gene's standardized variance that is sampling noise. The latent correlation
matrix is estimated as $R_x-\operatorname{diag}(\nu_jR_{x,jj})$, after the off-diagonal entries of $R_x$
have been multiplied by 0.9, with eigenvalues floored at 0.001 times the mean diagonal. The Spearman--Brown step treats the two
half-depth values as parallel measurements, which is exact on the count scale
and approximate after the logarithm, and library normalization couples genes;
in simulations with known latent covariance both effects were small
(below).

The proof of (3) shows that $\|F(\Sigma)\|^2_{\mathrm{op}}$ is the largest
eigenvalue of the channel applied to $(V^\top\Sigma^{-1}V)^{-1}$, the variance
of the $M$-coordinates of $x$ that the orthogonal coordinates leave unexplained.
Noise raises this conditional variance, so the statistic computed with measured
covariance bounds the latent one from above. After the noise diagonal is removed
from a covariance estimated from $n$ cells, fluctuations of the noise covariance
of spectral norm of order $\nu(p/n)^{1/2}$ remain off the diagonal\cite{marchenko1967}.
They create near-null directions that shrink the conditional variance, and in
simulations the statistic computed with the corrected covariance understated
the latent one, so that the two computed statistics bracket it
(below).

\subsection*{Validation of the noise correction}
Each of \RcBmSamples{} bone-marrow samples of the NeurIPS 2021 CITE-seq data\cite{luecken2021sandbox} (adaptation cells) defined a
latent truth. Each cell's expected proportions of the panel genes were the
average of the observed proportions (count divided by library) over the cell and
its 14 nearest neighbours in 30 principal components of standardized log
expression, and the rest of the transcriptome kept the remaining share. In a
second construction, independent gene-specific lognormal variation was added,
sized so that each gene's latent within-type variance matched the real data's
within-type variance times one minus its count-splitting noise share. Expected
libraries were the observed ones times 0.5, 1 or 2. Counts of the panel genes and
of the rest of the transcriptome were drawn from Poisson distributions, and the
library was their sum, so normalization coupled genes as in real data. Forty
independent replicate draws per cell gave the noise covariance
$N=E\{\operatorname{Cov}(x\mid\text{truth})\}$, a full matrix, and the latent
covariance $\operatorname{Cov}\{E(x\mid\text{truth})\}$, estimated as the
covariance of replicate means minus $N/40$, for
$x=\log(1+10^4\,\text{count}/\text{library})$. The first replicate was analysed
by the pipeline of the pairing-free analyses, and the truth was expressed on the
same scale (standard deviations of the analysed replicate, same shrinkage), with
the pairing-free channels of Supplementary Note~\ref{sec:cross} and the sample's real protein correlation.

Estimated noise shares had a median absolute error of
\RcSimNuMaeMin--\RcSimNuMaeMax{} and a median bias of at most \RcSimNuBiasMax{} in
every setting (Supplementary Table~\ref{tab:sim}). Off-diagonal noise covariance,
which normalization creates, was at most \RcSimCouplingMax\% of the off-diagonal
latent covariance in norm. Correcting the diagonal reduced the relative error of
the RNA correlation matrix from \RcSimErrMeasMin--\RcSimErrMeasMax{} to
\RcSimErrCorrMin--\RcSimErrCorrMax, as much as correction with the true noise
variances did. The compatibility statistic was nevertheless biased: across
settings and channels, its median ratio to the latent value was
\RcSimCorrRatioMin--\RcSimCorrRatioMax{} with corrected covariance
(\RcSimFrozenRatioMin--\RcSimFrozenRatioMax{} for the prespecified channel) and
\RcSimMeasRatioMin--\RcSimMeasRatioMax{} with measured covariance, and the latent
value lay between the two in \RcSimBracket{} of \RcSimBracketTotal{} cases. Correction
with the true noise variances gave nearly the same statistics (median ratio to
the estimated correction, \RcSimOracleMin--\RcSimOracleMax), so the bias arises
from finite-cell fluctuations of the noise covariance rather than from the
estimated shares (Proposition~\ref{prop:noise} and the remark after it). At the observed
depth, the prespecified blood channel's latent statistic exceeded one in all \RcSimDOne{}
simulated total analyses (median \RcSimFrozenTrue), while its corrected statistic
exceeded one in \RcSimFrozenCorrAbove.

\begin{table}[h]
\caption{\textbf{Validation of the noise correction.} Medians over
\RcBmSamples{} bone-marrow samples. Noise share, true median share of sampling
noise in gene variance; share error, mean absolute error of the count-splitting
estimate; matrix error, relative Frobenius error of the measured and corrected RNA
correlation matrices against the latent one (all with the same shrinkage);
$\|F\|$ / true, ratio of the compatibility statistic of the prespecified blood channel
computed with measured or corrected covariance to its latent value. Depth, factor
applied to the observed library sizes.}
\label{tab:sim}
\centering\footnotesize\setlength{\tabcolsep}{4pt}
% Generated by extension/recoverability/make_assets.py. Do not edit.
\begin{tabular}{llrcccccc}
\toprule
& & & & & \multicolumn{2}{c}{Matrix error} & \multicolumn{2}{c}{$\|F\|$ / true, blood channel} \\
\cmidrule(lr){6-7}\cmidrule(lr){8-9}
Latent & Cells & Depth & Noise share & Share error & Measured & Corrected & Measured & Corrected \\
\midrule
Smoothed & All & 0.5 & 0.56 & 0.055 & 0.29 & 0.11 & 2.25 & 0.23 \\
 & All & 1 & 0.39 & 0.041 & 0.18 & 0.08 & 1.76 & 0.27 \\
 & All & 2 & 0.25 & 0.028 & 0.12 & 0.06 & 1.46 & 0.47 \\
 & Within types & 0.5 & 0.78 & 0.067 & 0.95 & 0.31 & 2.79 & 0.22 \\
 & Within types & 1 & 0.64 & 0.057 & 0.61 & 0.23 & 2.18 & 0.23 \\
 & Within types & 2 & 0.49 & 0.041 & 0.38 & 0.16 & 1.79 & 0.26 \\
Plus gene-specific & All & 0.5 & 0.48 & 0.055 & 0.32 & 0.13 & 1.92 & 0.27 \\
 & All & 1 & 0.33 & 0.039 & 0.20 & 0.10 & 1.55 & 0.29 \\
 & All & 2 & 0.21 & 0.027 & 0.13 & 0.07 & 1.37 & 0.63 \\
 & Within types & 0.5 & 0.66 & 0.070 & 1.01 & 0.38 & 2.09 & 0.25 \\
 & Within types & 1 & 0.52 & 0.055 & 0.67 & 0.28 & 1.67 & 0.27 \\
 & Within types & 2 & 0.36 & 0.038 & 0.44 & 0.20 & 1.43 & 0.41 \\
\bottomrule
\end{tabular}

\end{table}
